\clearpage
\fancyhead[L]{\small DETERMINATION OF THE ACCURACY OF OBSERVATIONS}
\pdfbookmark[0]{Determination of the Accuracy of Observations}{genauigkeit}
\SourceLeaf{114}
\begin{center}
\rule{.8\linewidth}{.3pt}\par\smallskip
{\small Zeitschrift für Astronomie, edited by B. von \textsc{Lindenau} and J. G. F. \textsc{Bohnenberger}.\par
Volume I. Issue for March and April 1816.\par}
\smallskip\rule{.8\linewidth}{.3pt}\par\bigskip
{\large\itshape Determination of the Accuracy of Observations.\par}
\end{center}
\renewcommand{\Article}[1]{\par\medskip\Needspace{4\baselineskip}\pdfbookmark[1]{Art. #1}{genauigkeit-art-#1}\centerline{\textbf{#1.}}\nobreak\smallskip}
\Article{1}
In establishing the so-called method of least squares, it is assumed that the probability of an observational error \(\Delta\) is expressed by the formula
\[
\frac{h}{\surd\pi}.e^{-hh\Delta\Delta}
\]
where \(\pi\) denotes half the circumference of a circle, \(e\) the base of hyperbolic logarithms, and \(h\) a constant which, as in Art. 178 of the \textit{Theoria Motus Corporum Coelestium}, may be regarded as the measure of the accuracy of the observations. When the method of least squares is applied to determine the most probable values of the quantities on which the observations depend, there is no need to know the value of \(h\); the \textit{ratio} of the accuracy of the results to that of the observations is also independent of \(h\). Nevertheless, knowledge of this quantity itself is always interesting and instructive, and I shall therefore show how such knowledge may be obtained from the observations themselves.

\Article{2}
I shall first offer a few remarks to explain the subject. For brevity, I denote the value of the integral
\SourceLeaf{115}
\[
\int\frac{2e^{-tt}\,\mathrm{d}t}{\surd\pi}
\]
taken from \(t=0\) by \(\theta t\). A few individual values will give an idea of the course of this function. We have
\[
\begin{array}{r@{\;=\;}l@{\;=\;}l}
0{,}5000000&\theta\,0{,}4769363&\theta\rho\\
0{,}6000000&\theta\,0{,}5951161&\theta\,1{,}247790\rho\\
0{,}7000000&\theta\,0{,}7328691&\theta\,1{,}536618\rho\\
0{,}8000000&\theta\,0{,}9061939&\theta\,1{,}900032\rho\\
0{,}8427008&\theta\,1&\theta\,2{,}096716\rho\\
0{,}9000000&\theta\,1{,}1630872&\theta\,2{,}438664\rho\\
0{,}9900000&\theta\,1{,}8213864&\theta\,3{,}818930\rho\\
0{,}9990000&\theta\,2{,}3276754&\theta\,4{,}880475\rho\\
0{,}9999000&\theta\,2{,}7510654&\theta\,5{,}768204\rho\\
1&\multicolumn{2}{l}{\theta\infty}
\end{array}
\]

The probability that the error of an observation lies between the limits \(-\Delta\) and \(+\Delta\), or, irrespective of its sign, is no greater than \(\Delta\), is
\[
=\int\frac{he^{-hhxx}\,\mathrm{d}x}{\surd\pi}
\]
if the integral is taken from \(x=-\Delta\) to \(x=+\Delta\); or it is twice as great as the same integral taken from \(x=0\) to \(x=\Delta\), and hence
\[
=\theta h\Delta
\]

The probability that the error is not less than \(\frac\rho h\) is therefore \(=\frac12\), or equal to the probability of the contrary event. We shall call this quantity \(\frac\rho h\) the \textit{probable error}, and denote it by \(r\). By contrast, the probability that the error exceeds \(2{,}438664r\) is only \(\frac1{10}\); the probability that the error exceeds \(3{,}818930r\) is only \(\frac1{100}\), and so on.

\Article{3}
Let us now suppose that in \(m\) actual observations the errors \(\alpha,\,\beta,\,\gamma,\,\delta\), and so on, have been made, and examine what may be inferred from them about the values of \(h\) and \(r\). If we make two \SourceLeaf{116}%
assumptions, taking the true value of \(h\) to be either \(=H\) or \(=H'\), then the probabilities with which the errors \(\alpha,\,\beta,\,\gamma,\,\delta\), and so on, would be expected under them are respectively in the ratio

\[
\begin{gathered}
He^{-HH\alpha\alpha}.He^{-HH\beta\beta}.He^{-HH\gamma\gamma}.\text{ etc.}\\
\text{to}\quad H'e^{-H'H'\alpha\alpha}.H'e^{-H'H'\beta\beta}.H'e^{-H'H'\gamma\gamma}.\text{ etc.}
\end{gathered}
\]
that is, in the ratio
\[
H^m e^{-HH(\alpha\alpha+\beta\beta+\gamma\gamma+\text{ etc.})}\quad\text{to}\quad
H'{}^m e^{-H'H'(\alpha\alpha+\beta\beta+\gamma\gamma+\text{ etc.})}
\]
Consequently, \textit{after} the occurrence of those errors, the probabilities that \(H\) or \(H'\) was the true value of \(h\) stand in the same ratio (\textit{T. M. C. C.}, Art. 176); or the probability of each value of \(h\) is proportional to
\[
h^m e^{-hh(\alpha\alpha+\beta\beta+\gamma\gamma+\text{ etc.})}
\]
The \textit{most probable} value of \(h\) is therefore the one for which this quantity is a maximum, and, by the known rules, is found to be
\[
=\surd\frac{m}{2(\alpha\alpha+\beta\beta+\gamma\gamma+\text{ etc.})}
\]
The most probable value of \(r\) is consequently
\[
\begin{aligned}
&=\rho\surd\frac{2(\alpha\alpha+\beta\beta+\gamma\gamma+\text{ etc.})}{m}\\[4pt]
&=0{,}6744897\surd\frac{\alpha\alpha+\beta\beta+\gamma\gamma+\text{ etc.}}{m}
\end{aligned}
\]
This result is general, whether \(m\) is large or small.

\Article{4}
It is easy to see that, for this determination of \(h\) and \(r\), the smaller \(m\) is, the less reason one has to expect a high degree of accuracy. Let us therefore develop the degree of accuracy to be assigned to this determination for the case where \(m\) is a large number. For brevity, let us denote the probable value of \(h\) that has been found
\[
\surd\frac{m}{2(\alpha\alpha+\beta\beta+\gamma\gamma+\text{ etc.})}
\]
by \(H\), and note that the probability that \(H\) is the
\SourceLeaf{117}%
true value of \(h\) is to the probability that the true value is \(=H+\lambda\) as
\[
H^m e^{-\frac m2}:(H+\lambda)^m e^{-\frac{m(H+\lambda)^2}{2HH}}
\]
or as
\[
1:e^{-\frac{\lambda\lambda m}{HH}(1-\frac13.\frac\lambda H+\frac{1\,\lambda\lambda}{4HH}-\frac15.\frac{\lambda^3}{H^3}\, .\,.\,.\,.)}
\]
The second term will be appreciable compared with the first only when \(\frac\lambda H\) is a small fraction; hence we may allow ourselves to use the following in place of the ratio given:
\[
1:e^{-\frac{\lambda\lambda m}{HH}}
\]
This in fact means that the probability that the true value of \(h\) lies between \(H+\lambda\) and \(H+\lambda+\mathrm{d}\lambda\) is very nearly
\[
=Ke^{-\frac{\lambda\lambda m}{HH}}\mathrm{d}\lambda
\]
where \(K\) is a constant that must be determined so that the integral
\[
\int Ke^{-\frac{\lambda\lambda m}{HH}}\mathrm{d}\lambda
\]
taken between the admissible limits of \(\lambda\), becomes \(=1\). Here, instead of such limits, since owing to the largeness of \(m\),
\[
e^{-\frac{\lambda\lambda m}{HH}}
\]
evidently becomes imperceptible as soon as \(\frac\lambda H\) ceases to be a small fraction, it is permissible to take the limits \(-\infty\) and \(+\infty\), which gives
\[
K=\frac1H\surd\frac m\pi
\]
Hence the probability that the true value of \(h\) lies between \(H-\lambda\) and \(H+\lambda\) is
\[
=\theta\bigl(\tfrac\lambda H\surd m\bigr)
\]
so that this probability is \(=\frac12\) when
\[
\frac\lambda H\surd m=\rho\quad\text{is.}
\]
\SourceLeaf{118}%
The odds are even that the true value of \(h\) lies
\[
\text{between}\quad H\bigl(1-\tfrac\rho{\surd m}\bigr)\quad\text{and}\quad H\bigl(1+\tfrac\rho{\surd m}\bigr)
\]
or that the true value of \(r\) lies
\[
\text{between}\quad\frac{R}{1-\frac\rho{\surd m}}\quad\text{and}\quad\frac{R}{1+\frac\rho{\surd m}}
\]
where \(R\) denotes the most probable value of \(r\) found in the preceding article. These limits may be called the \textit{probable limits of the true values of \(h\) and \(r\)}; evidently, for the probable limits of the true value of \(r\), we may also put
\[
R\bigl(1-\tfrac\rho{\surd m}\bigr)\quad\text{and}\quad R\bigl(1+\tfrac\rho{\surd m}\bigr)
\]

\Article{5}
In the preceding investigation, we took the view that \(\alpha,\,\beta,\,\gamma,\,\delta\), and so on, were definite and given quantities, and sought the probability that the true value of \(h\) or \(r\) lies between certain limits. The question can also be considered from another point of view: assuming that the observation errors obey some definite law of probability, we can determine the probability that the sum of the squares of \(m\) observation errors falls between certain limits. For large \(m\), this problem has already been solved by \textsc{Laplace}, as has the problem of finding the probability that the sum of \(m\) observation errors itself falls between certain limits. The investigation can readily be generalized further; I shall give only the result here.

Let \(\varphi x\) denote the probability of the observation error \(x\), so that \(\int\varphi x.\mathrm{d}x=1\) when the integral is taken from \(x=-\infty\) to \(x=+\infty\). Between the same limits, let us generally denote the value of the integral
\[
\int\varphi x.x^n\mathrm{d}x
\]
\SourceLeaf{119}
by \(K^{\mathrm n}\).\footnote{Or rather, the integral \(\int\varphi x.x^n\mathrm{d}x\) from \(x=0\) to \(x=\infty\) is to be denoted by
\[\tfrac12K^{\mathrm n}\]
[Handwritten remark].} Let \(S^{\mathrm n}\) be the sum
\[
\alpha^n+\beta^n+\gamma^n+\delta^n+\text{ etc.}
\]
where \(\alpha,\,\beta,\,\gamma,\,\delta\), and so on, represent any \(m\) observation errors; all terms in this sum are to be taken as positive, even when \(n\) is odd.

Then \(mK^{\mathrm n}\) is the most probable value of \(S^{\mathrm n}\), and the probability that the true value of \(S^{\mathrm n}\) lies between the limits \(mK^{\mathrm n}-\lambda\) and \(mK^{\mathrm n}+\lambda\) is
\[
=\theta\frac{\lambda}{\surd[2m(K^{2\mathrm n}-K^{\mathrm n}K^{\mathrm n})]}
\]
Consequently, the probable limits of \(S^{\mathrm n}\) are
\[
mK^{\mathrm n}-\rho\surd[2m(K^{2\mathrm n}-K^{\mathrm n}K^{\mathrm n})]\quad\text{and}\quad
mK^{\mathrm n}+\rho\surd[2m(K^{2\mathrm n}-K^{\mathrm n}K^{\mathrm n})]
\]
This result holds generally for every law of observation errors. Apply it to the case in which
\[
\varphi x=\frac h{\surd\pi}e^{-hhxx}
\]
is set; we then obtain
\[
K^{\mathrm n}=\frac{\Pi\,\frac12(n-1)}{h^n\surd\pi}
\]
with the characteristic \(\Pi\) taken in the sense of the \textit{Disquisitiones generales circa seriem infinitam} (Comm. nov. soc. Gotting. T. II.), M. 5, Art. 28 of the cited treatise. Thus
\[
\begin{aligned}
K&=1,&K'&=\frac1{h\surd\pi},&K''&=\frac1{2hh},&K'''&=\frac1{h^3\surd\pi}\\[5pt]
K^{\mathrm{IV}}&=\frac{1.3}{4h^4},&K^{\mathrm{V}}&=\frac{1.2}{h^5\surd\pi},&K^{\mathrm{VI}}&=\frac{1.3.5}{8h^6},&K^{\mathrm{VII}}&=\frac{1.2.3}{h^7\surd\pi}\text{ etc.}
\end{aligned}
\]
Consequently, the most probable value of \(S^{\mathrm n}\) is
\[
\frac{m\Pi\,\frac12(n-1)}{h^n\surd\pi}
\]
and the probable limits of the true value of \(S^{\mathrm n}\) are
\SourceLeaf{120}
\[
\frac{m\Pi\,\frac12(n-1)}{h^n\surd\pi}\Bigl\{1-\rho\surd\Bigl(\frac2m.\Bigl(\frac{\Pi(n-\frac12).\surd\pi}{(\Pi\,\frac12(n-1))^2}-1\Bigr)\Bigr)\Bigr\}
\]
and
\[
\frac{m\Pi\,\frac12(n-1)}{h^n\surd\pi}\Bigl\{1+\rho\surd\Bigl(\frac2m.\Bigl(\frac{\Pi(n-\frac12).\surd\pi}{(\Pi\,\frac12(n-1))^2}-1\Bigr)\Bigr)\Bigr\}
\]
If, as above, we set
\[
\frac\rho h=r
\]
so that \(r\) represents the probable observation error, then the most probable value of
\[
\rho\BareRoot[n]{\frac{S^{\mathrm n}\surd\pi}{m\Pi\,\frac12(n-1)}}
\]
is evidently \(=r\); and the probable limits of the value of that quantity are
\[
r\Bigl\{1-\frac\rho n\surd\Bigl(\frac2m.\Bigl(\frac{\Pi(n-\frac12).\surd\pi}{(\Pi\,\frac12(n-1))^2}-1\Bigr)\Bigr)\Bigr\}
\]
and
\[
r\Bigl\{1+\frac\rho n\surd\Bigl(\frac2m.\Bigl(\frac{\Pi(n-\frac12).\surd\pi}{(\Pi\,\frac12(n-1))^2}-1\Bigr)\Bigr)\Bigr\}
\]
Thus one may also wager even odds that \(r\) lies between the limits
\[
\begin{aligned}
&\rho\BareRoot[n]{\frac{S^{\mathrm n}\surd\pi}{m\Pi\,\frac12(n-1)}}\\[3pt]
&\quad\Bigl\{1-\frac\rho n\surd\Bigl(\frac2m.\Bigl(\frac{\Pi(n-\frac12).\surd\pi}{(\Pi\,\frac12(n-1))^2}-1\Bigr)\Bigr)\Bigr\}
\end{aligned}
\]
and
\[
\begin{aligned}
&\rho\BareRoot[n]{\frac{S^{\mathrm n}\surd\pi}{m\Pi\,\frac12(n-1)}}\\[3pt]
&\quad\Bigl\{1+\frac\rho n\surd\Bigl(\frac2m.\Bigl(\frac{\Pi(n-\frac12).\surd\pi}{(\Pi\,\frac12(n-1))^2}-1\Bigr)\Bigr)\Bigr\}
\end{aligned}
\]
For \(n=2\), these limits are
\[
\rho\surd\frac{2S''}{m}\bigl\{1-\tfrac\rho{\surd m}\bigr\}\quad\text{and}\quad
\rho\surd\frac{2S''}{m}\bigl\{1+\tfrac\rho{\surd m}\bigr\}
\]
in full agreement with those found above (Art. 4). In general, for even \(n\), the limits are
\[
\begin{aligned}
&\rho\surd2.\BareRoot[n]{\frac{S^{\mathrm n}}{m.1.3.5.7\, .\,.\,.\,.\,(n-1)}}\\[4pt]
&\quad\Bigl\{1-\frac\rho n\surd\Bigl(\frac2m.\Bigl(\frac{(n+1).(n+3)\, .\,.\,.\,(2n-1)}{1.3.5\, .\,.\,.\,(n-1)}-1\Bigr)\Bigr)\Bigr\}
\end{aligned}
\]
and
\[
\begin{aligned}
&\rho\surd2.\BareRoot[n]{\frac{S^{\mathrm n}}{m.1.3.5.7\, .\,.\,.\,.\,(n-1)}}\\[4pt]
&\quad\Bigl\{1+\frac\rho n\surd\Bigl(\frac2m.\Bigl(\frac{(n+1).(n+3)\, .\,.\,.\,(2n-1)}{1.3.5\, .\,.\,.\,(n-1)}-1\Bigr)\Bigr)\Bigr\}
\end{aligned}
\]
and for odd \(n\), the following:
\[
\begin{aligned}
&\rho\BareRoot[n]{\frac{S^{\mathrm n}\surd\pi}{m.1.2.3\, .\,.\,.\,\frac12(n-1)}}\\[4pt]
&\quad\Bigl\{1-\frac\rho n\surd\Bigl(\frac1m.\Bigl(\frac{1.3.5.7\, .\,.\,.\,(2n-1)\pi}{(2.4.6\, .\,.\,.\,(n-1))^2}-2\Bigr)\Bigr)\Bigr\}
\end{aligned}
\]
and
\[
\begin{aligned}
&\rho\BareRoot[n]{\frac{S^{\mathrm n}\surd\pi}{m.1.2.3\, .\,.\,.\,\frac12(n-1)}}\\[4pt]
&\quad\Bigl\{1+\frac\rho n\surd\Bigl(\frac1m.\Bigl(\frac{1.3.5.7\, .\,.\,.\,(2n-1)\pi}{(2.4.6\, .\,.\,.\,(n-1))^2}-2\Bigr)\Bigr)\Bigr\}
\end{aligned}
\]

\SourceLeaf{121}
\Article{6}
I add the numerical values for the simplest cases:
\begin{center}Probable limits of \(r\)\end{center}
\[
\begin{array}{r@{\quad}l}
\mathrm{I}.&0{,}8453473\quad\dfrac{S'}m.\bigl(1\mp\dfrac{0{,}5095841}{\surd m}\bigr)\\[10pt]
\mathrm{II}.&0{,}6744897\BareRoot[2]{\dfrac{S''}m}.\bigl(1\mp\dfrac{0{,}4769363}{\surd m}\bigr)\\[10pt]
\mathrm{III}.&0{,}5771897\BareRoot[3]{\dfrac{S'''}m}.\bigl(1\mp\dfrac{0{,}4971987}{\surd m}\bigr)\\[10pt]
\mathrm{IV}.&0{,}5125017\BareRoot[4]{\dfrac{S''''}m}.\bigl(1\mp\dfrac{0{,}5507186}{\surd m}\bigr)\\[10pt]
\mathrm{V}.&0{,}4655532\BareRoot[5]{\dfrac{S^{\mathrm V}}m}.\bigl(1\mp\dfrac{0{,}6355080}{\surd m}\bigr)\\[10pt]
\mathrm{VI}.&0{,}4294972\BareRoot[6]{\dfrac{S^{\mathrm{VI}}}m}.\bigl(1\mp\dfrac{0{,}7557764}{\surd m}\bigr)
\end{array}
\]
This also shows that method II is the most advantageous of them all. One hundred observation errors treated by this formula give a result just as reliable as
\begin{center}
114 under I,\quad109 under III,\quad133 under IV,\quad178 under V,\quad251 under VI.
\end{center}
Formula I, however, has the advantage of being by far the easiest to calculate. Since it is not much less accurate than II, it may certainly be used when one neither already knows the sum of the squares of the errors for other reasons nor wishes to know it.

\Article{7}
An easier, though considerably less accurate, procedure is the following. Arrange all \(m\) observation errors in order of their absolute values, and call the middle one \(M\) when their number is odd, or the arithmetic mean of the two middle ones when their number is even. It can be shown—but this cannot be worked out further here—that for a large number of observations, \(r\) is the most probable value of \(M\), and that the probable limits of \(M\) are
\[
r\bigl(1-e^{\rho\rho}\surd\tfrac\pi{8m}\bigr)\quad\text{and}\quad r\bigl(1+e^{\rho\rho}\surd\tfrac\pi{8m}\bigr)
\]
or the probable limits of the value of \(r\) are
\SourceLeaf{122}
\[
\begin{gathered}
M\bigl(1-e^{\rho\rho}\surd\tfrac\pi{8m}\bigr)\quad\text{and}\quad M\bigl(1+e^{\rho\rho}\surd\tfrac\pi{8m}\bigr),\\[4pt]
\text{or numerically}\quad M\bigl(1\mp\tfrac{0{,}7520974}{\surd m}\bigr)
\end{gathered}
\]
This procedure is therefore only a little more accurate than applying formula VI; one would have to consult 249 observation errors to go as far as one does with 100 observation errors under formula II.

\Article{8}
Applying some of these methods to the errors reported in \textsc{Bode}'s Astronomical Yearbook for 1818, p. 234, from 48 observations of the right ascensions of the Pole Star by \textsc{Bessel}, gives
\[
S'=60''46,\quad S''=110''600,\quad S'''=250''341118
\]
The most probable values of \(r\) that follow are
\begin{center}
\begin{tabular}{r@{\;}c@{\;}l@{\quad}l@{\;}l}
by formula I&.\,.\,.\,.&\(1''065,\)&probable uncertainty \(=\)&\(\pm0''078\)\\
II&.\,.\,.\,.&\(1{,}024\)&&\(\pm0{,}070\)\\
III&.\,.\,.\,.&\(1{,}001\)&&\(\pm0{,}072\)\\
by Art. 7&.\,.\,.\,.\,.&\(1{,}045\)&&\(\pm0{,}113\)
\end{tabular}
\end{center}
a degree of agreement hardly to be expected. \textsc{Bessel} himself gives \(1''067\), and therefore seems to have calculated according to formula I.
\begin{center}\rule{.24\linewidth}{.3pt}\end{center}
